% !TeX program = pdflatex
\documentclass[aspectratio=169,xcolor={dvipsnames,table}]{beamer}

\input{../../_en_preambulo_beamer}
% Comandos y macros estadísticas compartidas para las presentaciones Beamer en inglés (Metropolis)

% Conjuntos numéricos básicos
\newcommand{\R}{\mathbb{R}}
\newcommand{\N}{\mathbb{N}}
\newcommand{\Q}{\mathbb{Q}}
\newcommand{\Z}{\mathbb{Z}}
\newcommand{\set}[1]{\left\{ #1 \right\}}
\newcommand{\abs}[1]{\left|#1\right|}
\newcommand{\norm}[1]{\left\| #1 \right\|}

% Comandos estadísticos y probabilísticos del curso
\newcommand{\Var}{\operatorname{Var}}
\newcommand{\cov}{\operatorname{Cov}}
\newcommand{\corr}{\rho}
\newcommand{\comb}[2]{\begin{pmatrix} #1 \\ #2 \end{pmatrix}}
\newcommand{\s}{\sigma}
\newcommand{\particion}{\mathfrak{P}}
\newcommand{\card}[1]{\#\left( #1 \right)}
\newcommand{\seq}[3]{\set{#1}_{#2}^{#3}}
\newcommand{\E}{\mathbb{E}}
\newcommand{\Prob}{\mathbb{P}}

% Entornos pedagógicos y cajas en inglés académico natural (soportando tanto nombres en inglés como en español)
\newenvironment{discussion}{%
  \begin{alertblock}{Discussion Question}%
}{%
  \end{alertblock}%
}
\newenvironment{discusion}{%
  \begin{alertblock}{Discussion Question}%
}{%
  \end{alertblock}%
}

\newenvironment{reflection}{%
  \begin{alertblock}{Classroom Reflection}%
}{%
  \end{alertblock}%
}
\newenvironment{reflexion}{%
  \begin{alertblock}{Classroom Reflection}%
}{%
  \end{alertblock}%
}

\newenvironment{paradox}{%
  \begin{alertblock}{Exploratory Paradox}%
}{%
  \end{alertblock}%
}
\newenvironment{paradoja}{%
  \begin{alertblock}{Exploratory Paradox}%
}{%
  \end{alertblock}%
}

\newenvironment{motivation}{%
  \begin{exampleblock}{Motivation: Data Science in Practice}%
}{%
  \end{exampleblock}%
}
\newenvironment{motivacion}{%
  \begin{exampleblock}{Motivation: Data Science in Practice}%
}{%
  \end{exampleblock}%
}

\newenvironment{quickchallenge}{%
  \begin{block}{Quick Team Challenge}%
}{%
  \end{block}%
}
\newenvironment{retoclase}{%
  \begin{block}{Quick Team Challenge}%
}{%
  \end{block}%
}


\title{Continuous CDF and Quantiles}
\subtitle{Section 04.02 --- Cumulative Distribution, Inversion, and Monte Carlo}
\author[J. Castillo Colmenares]{Juliho Castillo Colmenares}
\institute[Tec de Monterrey]{Tecnologico de Monterrey}
\date{\vspace{-1.5cm}}

\begin{document}

% SLIDE 01: Title Page
\begin{frame}[plain]
	\titlepage
\end{frame}

% SLIDE 02: Roadmap
\begin{frame}{Roadmap --- Chapter 04: Continuous Random Variables}
	\vspace{-0.15cm}
	\begin{block}{Curricular Structure of Unit 3}
		\footnotesize
		\begin{enumerate}\itemsep=0.03cm
			\item \textcolor{gray}{Section 04.01: PDF and Continuous Support}
			\item \textbf{\textcolor{TecRojo}{Section 04.02: Continuous CDF and Quantiles}} $\leftarrow$ \emph{Current focus}
			\item \textcolor{gray}{Section 04.03: Expectation, Variance, and Continuous LOTUS}
			\item \textcolor{gray}{Section 04.04: Continuous Uniform Distribution}
			\item \textcolor{gray}{Section 04.05: Exponential Distribution and Memoryless Processes}
			\item \textcolor{gray}{Section 04.06: Normal Distribution and Z-Score}
			\item \textcolor{gray}{Section 04.07: Gamma, Beta, and Weibull Distributions}
		\end{enumerate}
	\end{block}
	\vspace{-0.2cm}
	\begin{alertblock}{Learning Objective}
		\scriptsize
		Construct the continuous Cumulative Distribution Function $F(x) = P(X \le x)$ from the PDF, master its axiomatic properties, compute probabilities via CDF differences, define quantiles as $F^{-1}(p)$, and apply the inverse transform method for random number generation.
	\end{alertblock}
\end{frame}

% SLIDE 03: Formal Definition
\begin{frame}{Formal Definition: CDF and Inverse Transform}
	\vspace{-0.15cm}
	\begin{columns}[T]
		\begin{column}{0.48\textwidth}
			\begin{block}{Continuous CDF}
				\footnotesize
				$F(x) = P(X \le x) = \int_{-\infty}^{x} f(t)\,dt$. By the Fundamental Theorem of Calculus, $f(x) = F'(x)$. The CDF summarizes the entire distributional information.
			\end{block}
		\end{column}
		\begin{column}{0.48\textwidth}
			\begin{alertblock}{Inverse Transform}
				\footnotesize
				If $U \sim U(0, 1)$ and $F$ is continuous and strictly increasing, then $X = F^{-1}(U)$ has CDF $F$. This enables exact sampling from any continuous distribution given only its CDF.
			\end{alertblock}
		\end{column}
	\end{columns}
\end{frame}

% SLIDE 04: Quantiles and KS Test
\begin{frame}{Quantiles and Kolmogorov-Smirnov Goodness-of-Fit}
	\vspace{-0.15cm}
	\begin{columns}[T]
		\begin{column}{0.48\textwidth}
			\begin{block}{Quantiles and Interval Probabilities}
				\footnotesize
				Quantile: $q_{p} = F^{-1}(p) = \inf\{x : F(x) \ge p\}$. Interval probability: $P(a \le X \le b) = F(b) - F(a)$. For Uniform(0,1): $F(x) = x$, $q_{p} = p$.
			\end{block}
		\end{column}
		\begin{column}{0.48\textwidth}
			\begin{alertblock}{Kolmogorov-Smirnov Test}
				\footnotesize
				$D_{n} = \sup_{x} |F_{n}(x) - F_{0}(x)|$. Under $H_{0}$: $\sqrt{n} D_{n} \xrightarrow{d} K$ (Kolmogorov distribution). Reject $H_{0}$ if $D_{n} > K_{1-\alpha}/\sqrt{n}$.
			\end{alertblock}
		\end{column}
	\end{columns}
\end{frame}

% SLIDE 05: Applications
\begin{frame}{Applications in Inference and Simulation}
	\vspace{-0.1cm}
	\begin{itemize}\itemsep=0.1cm
		\item \textbf{Confidence Intervals:} $\text{CI}_{95\%} = [\hat{\theta} \pm z_{0.025} \cdot \text{SE}(\hat{\theta})]$ uses $q_{0.025} = -1.96$ and $q_{0.975} = 1.96$. \pause
		\item \textbf{Monte Carlo Simulation:} Random variate generation via inversion for uncertainty propagation. \pause
		\item \textbf{Non-parametric Tests:} Kolmogorov-Smirnov and Anderson-Darling for model fit assessment. \pause
		\item \textbf{Risk Analysis:} Value-at-Risk (VaR) and Conditional VaR defined as distribution quantiles.
	\end{itemize}
\end{frame}

% SLIDE 06: Python Lab Block 1
\begin{frame}[fragile]{Python Lab: CDF Validation and Quantile Inversion}
	\vspace{-0.05cm}
	\scriptsize
	Validation of axiomatic CDF properties and quantile computation via numerical inversion (\texttt{04.02\_continuous\_cdf.py}):
	\vspace{0.02cm}
	{\lstset{basicstyle=\fontsize{4.5pt}{5.5pt}\selectfont\ttfamily}
	\lstinputlisting[firstline=18, lastline=43]{../../code/04_variables_aleatorias_continuas/04.02_continuous_cdf.py}}
\end{frame}

% SLIDE 07: Python Lab Block 2
\begin{frame}[fragile]{Python Lab: Inverse Transform Sampling}
	\vspace{-0.05cm}
	\scriptsize
	Implementation of $X = F^{-1}(U)$ for Exponential and Uniform with `scipy.stats.kstest`:
	\vspace{0.02cm}
	{\lstset{basicstyle=\fontsize{4pt}{5pt}\selectfont\ttfamily}
	\lstinputlisting[firstline=50, lastline=75]{../../code/04_variables_aleatorias_continuas/04.02_continuous_cdf.py}}
\end{frame}

% SLIDE 08: Python Lab Block 3
\begin{frame}[fragile]{Python Lab: KS Test and Empirical Quantiles}
	\vspace{-0.05cm}
	\scriptsize
	Application of Kolmogorov-Smirnov and comparison of empirical vs. theoretical quantiles:
	\vspace{0.02cm}
	{\lstset{basicstyle=\fontsize{4.5pt}{5.5pt}\selectfont\ttfamily}
	\lstinputlisting[firstline=85, lastline=110]{../../code/04_variables_aleatorias_continuas/04.02_continuous_cdf.py}}
\end{frame}

% SLIDE 09: Python Lab Terminal Output
\begin{frame}[fragile]{Python Lab: Terminal Output}
	\vspace{-0.1cm}
	\begin{block}{}
		\fontsize{4pt}{5pt}\selectfont
		\begin{verbatim}
=== Block 1: CDF Validation & Quantile Inversion ===
CDF boundary properties:
  Uniform(0,1): F(-inf)=0 | F(+inf)=1.000000
  Exponential(2.0): F(-inf)=0 | F(+inf)=1.000000
  Normal(0,1): F(-inf)=0 | F(+inf)=1.000000
Monotonicity: all True
Quantile inversion: all match=True

=== Block 2: Inverse Transform Sampling ===
Exponential(lambda=2):
  Empirical mean: 0.4980 (theoretical: 0.5000)
  Empirical var:  0.2465 (theoretical: 0.2500)
  KS test: stat=0.001974, p-value=0.8299
Uniform(0, 5):
  Empirical mean: 2.5067 (theoretical: 2.5000)
  KS test: stat=0.002729, p-value=0.4450

=== Block 3: Quantile-Based Statistics & KS Test ===
KS statistic for n=50 Uniform(0,1) samples: D=0.1420
  Critical value (alpha=0.05): 0.1921
  SciPy KS p-value: 0.2417
  Decision: Fail to reject H0
95% CI for median: [0.4231, 0.6110]
		\end{verbatim}
	\end{block}
\end{frame}

% SLIDE 10: Exercise Level 1 (Statement)
\begin{frame}{In-Class Exercise --- Level 1: Fundamental (1/2)}
	\vspace{-0.1cm}
	\begin{block}{Problem 4.2.2 --- Exponential CDF and Quantiles (Statement)}
		Consider the exponential distribution with parameter $\lambda = 2$ and PDF $f(x) = 2e^{-2x}$ for $x \ge 0$. Derive the CDF $F(x) = 1 - e^{-2x}$, calculate the median $m$ such that $F(m) = 0.5$, and find the quantiles $q_{0.25}$ and $q_{0.90}$.
	\end{block}
	\vspace{0.15cm}
	\pause
	\begin{alertblock}{Model Setup}
		\begin{itemize}\itemsep=0.04cm
			\item Support: $[0, \infty)$.
			\item CDF: $F(x) = \int_{0}^{x} 2e^{-2t}\,dt = 1 - e^{-2x}$ for $x \ge 0$.
			\item Quantile: $q_{p}$ such that $1 - e^{-2q_{p}} = p$, giving $q_{p} = -\frac{1}{2}\ln(1-p)$.
		\end{itemize}
	\end{alertblock}
\end{frame}

% SLIDE 11: Exercise Level 1 (Solution)
\begin{frame}{In-Class Exercise --- Level 1: Fundamental (2/2)}
	\vspace{-0.05cm}
	\scriptsize
	Solve step by step: \pause
	\begin{block}{Median and Quantiles}
		\begin{itemize}\itemsep=0.05cm
			\item \textbf{Median} $m$: $1 - e^{-2m} = 0.5$ implies $e^{-2m} = 0.5$, so $m = \frac{1}{2}\ln 2 \approx 0.347$. \pause
			\item \textbf{Percentile 25} $q_{0.25}$: $q_{0.25} = -\frac{1}{2}\ln(0.75) \approx 0.144$. \pause
			\item \textbf{Percentile 90} $q_{0.90}$: $q_{0.90} = -\frac{1}{2}\ln(0.10) \approx 1.151$.
		\end{itemize}
	\end{block}
	\vspace{0.05cm}
	\pause
	\begin{alertblock}{Interpretation}
		\scriptsize
		$25\%$ of observations are $\le 0.144$, $50\%$ are $\le 0.347$, and $90\%$ are $\le 1.151$. Quantiles grow logarithmically: extending to thicker tails requires many more samples.
	\end{alertblock}
\end{frame}

% SLIDE 12: Exercise Level 2 (Statement)
\begin{frame}{In-Class Exercise --- Level 2: Operational (1/2)}
	\vspace{-0.1cm}
	\begin{block}{Problem 4.2.5 --- Bank Waiting Time (Statement)}
		The waiting time in a bank queue (in minutes) follows $X \sim \text{Exp}(\lambda = 0.2)$. Calculate $P(X \le 3)$ and $P(X \ge 10)$, find the quantile $q_{0.90}$, and determine $t$ such that $P(X \le t) = 0.95$.
	\end{block}
	\vspace{0.15cm}
	\pause
	\begin{alertblock}{Practical Application}
		\scriptsize
		The bank wants to ensure that 95\% of customers wait less than $t$ minutes. This question is fundamental for designing the customer experience and optimizing staff levels.
	\end{alertblock}
\end{frame}

% SLIDE 13: Exercise Level 2 (Solution)
\begin{frame}{In-Class Exercise --- Level 2: Operational (2/2)}
	\vspace{-0.05cm}
	\scriptsize
	Apply the exponential CDF $F(x) = 1 - e^{-0.2 x}$: \pause
	\begin{block}{Probabilities and Quantiles}
		\begin{itemize}\itemsep=0.05cm
			\item $P(X \le 3) = 1 - e^{-0.6} \approx 0.4512$. \pause
			\item $P(X \ge 10) = e^{-2} \approx 0.1353$. \pause
			\item $q_{0.90}$: $e^{-0.2 q} = 0.1$ implies $q_{0.90} = -5\ln(0.1) = 5\ln 10 \approx 11.51$ min. \pause
			\item $t$ for 95\%: $e^{-0.2 t} = 0.05$ implies $t = -5\ln(0.05) \approx 14.98$ min.
		\end{itemize}
	\end{block}
	\vspace{0.05cm}
	\pause
	\begin{alertblock}{Management Decision}
		\scriptsize
		To serve 95\% of customers in $\le 15$ min, the bank needs to size the system with $\lambda \cdot t \ge 3$, i.e., enough service capacity to avoid long queues during peak hours.
	\end{alertblock}
\end{frame}

% SLIDE 14: Exercise Level 3 (Statement)
\begin{frame}{In-Class Exercise --- Level 3: Analytical (1/2)}
	\vspace{-0.1cm}
	\begin{block}{Problem 4.2.7 --- CDF Continuity (Statement)}
		Prove formally that the CDF $F(x) = P(X \le x)$ of a continuous random variable is a continuous function on all $\mathbb{R}$. Use the property $P(X = x) = 0$ and the law of total probability.
	\end{block}
	\vspace{0.1cm}
	\pause
	\begin{alertblock}{Deduction Strategy}
		\scriptsize
		For $h > 0$, use $|F(x+h) - F(x)| = |P(x < X \le x+h)|$ and the convergence of infinitesimal interval probabilities to 0.
	\end{alertblock}
\end{frame}

% SLIDE 15: Exercise Level 3 (Solution)
\begin{frame}{In-Class Exercise --- Level 3: Analytical (2/2)}
	\vspace{-0.05cm}
	\scriptsize
	Demonstrate right-continuity: \pause
	\begin{block}{Continuity Proof}
		\scriptsize
		For $h > 0$:
		\begin{align*}
			0 \le F(x + h) - F(x) = P(x < X \le x + h) \le P(x - h < X \le x + h) = F(x+h) - F(x-h).
		\end{align*}
		Since $X$ is continuous, $P(X = x) = 0$. By monotonicity of probability:
		\begin{align*}
			\lim_{h \to 0^{+}} P(x < X \le x + h) = P(\emptyset) = 0.
		\end{align*}
		Therefore $\lim_{h \to 0^{+}} F(x + h) = F(x)$. Left-continuity follows analogously. \quad \qedhere
	\end{block}
\end{frame}

% SLIDE 16: Exercise Level 4 (Statement)
\begin{frame}{In-Class Exercise --- Level 4: Challenging (1/2)}
	\vspace{-0.1cm}
	\begin{block}{Problem 4.2.10 --- Kolmogorov-Smirnov Test (Statement)}
		Under $H_{0}$ (data from $F_{0}$), prove that $\sqrt{n} D_{n} \xrightarrow{d} K$ where $K$ is the Kolmogorov distribution. For $n = 50$ Uniform(0, 1) samples with $D_{50} = 0.18$, determine whether $H_{0}$ is rejected at $\alpha = 0.05$.
	\end{block}
	\vspace{0.15cm}
	\pause
	\begin{alertblock}{Application}
		\scriptsize
		The critical value $K_{0.95} \approx 1.358$. Compare $\sqrt{50} \cdot 0.18$ with this threshold and conclude whether the data is consistent with Uniform(0, 1).
	\end{alertblock}
\end{frame}

% SLIDE 17: Exercise Level 4 (Solution)
\begin{frame}{In-Class Exercise --- Level 4: Challenging (2/2)}
	\vspace{-0.05cm}
	\scriptsize
	Apply the asymptotic Kolmogorov distribution: \pause
	\begin{block}{Test Statistic and Decision}
		\scriptsize
		For $n = 50$ and $D_{50} = 0.18$: $\sqrt{50} \cdot 0.18 \approx 1.273 < 1.358 = K_{0.95}$. We do not reject $H_{0}$ at $\alpha = 0.05$. \quad \qedhere
	\end{block}
	\vspace{0.05cm}
	\pause
	\begin{alertblock}{Interpretation and Application}
		\scriptsize
		The KS test is non-parametric: it does not require assuming a specific form for the alternative distribution. In practice, it is powerful for detecting deviations in the tail (where information is scarce) or median. For large $n$, the test is very sensitive to small deviations; hence Anderson-Darling is also used, which weights the tails more.
	\end{alertblock}
\end{frame}

% SLIDE 18: Closing and Next Steps
\begin{frame}{Closing Section and Modular Perspective}
	\vspace{-0.2cm}
	\begin{block}{Synthesis: The CDF as Inference Tool}
		\scriptsize
		The continuous CDF $F(x) = P(X \le x)$ encodes the entire distributional information and is the central tool for computing accumulated probabilities and defining quantiles. Its inverse $F^{-1}$ is the basis of percentile-based inference and the inverse transform method for Monte Carlo generation.
	\end{block}
	\vspace{0.05cm}
	\pause
	\begin{alertblock}{Key Lessons and Next Step}
		\begin{itemize}\itemsep=0.05cm
			\item \textbf{CDF vs. PDF:} $F$ is cumulative; $f$ is local. $f = F'$ and $F = \int f$.
			\item \textbf{Quantiles:} $q_{p} = F^{-1}(p)$ is the basis of confidence intervals and non-parametric tests.
			\item \textbf{Inverse Transform:} $X = F^{-1}(U)$ with $U \sim U(0, 1)$ generates exact samples from any continuous distribution.
			\item \textbf{Next Section:} \textcolor{TecRojo}{Expectation, Variance, and Continuous LOTUS}.
		\end{itemize}
	\end{alertblock}
\end{frame}

\end{document}
